\documentclass[11pt]{article}
\usepackage[utf8]{inputenc}
\usepackage{amsmath,amssymb,amsthm}
\usepackage[margin=1in]{geometry}
\usepackage{hyperref}
\usepackage{xcolor}
\usepackage{enumitem}
\usepackage{newunicodechar}
\newunicodechar{ℝ}{\ensuremath{\mathbb{R}}}
\newunicodechar{ℕ}{\ensuremath{\mathbb{N}}}

\title{Tide retrospective:\\\emph{2D partition function asymptotic (Z2)}}
\author{Daniel Murfet}
\date{18 May 2026}

\begin{document}
\maketitle

\begin{abstract}
Completes today's V2/Z1/Z2 triple --- moment, numerator, partition,
all packaged as exact reformulation + rescaled \texttt{Tendsto} +
\texttt{Asymptotics.IsEquivalent} for the 2D quartic-sextic potential.
Z2 reduces to Z1 specialised to $(j, k) = (0, 0)$: the seabed already
supplies the exact closed form for the 2D partition function, so
this Tide is purely the asymptotic-packaging step. 125 lines,
single-attempt clean build, no roadblocks, no GPT consult (third
consecutive mirror-shaped tide today). Zero \texttt{sorry}, zero
\texttt{axiom}, zero \texttt{native\_decide}.
\end{abstract}

\section{Setting}

V2 packaged T1's exact moment closed form; Z1 packaged the
unnormalised numerator's exact closed form. The partition function
is the natural third member of the triple: $Z_{2D}(t) =
\int\!\int e^{-tL(x, y)}\,dx\,dy$ for $L(x, y) = x^4/24 + y^6/720$.
Its exact closed form was already in the seabed, named
\texttt{partitionFunction\_quarticSextic\_eq}; Z2 just adds the
asymptotic packaging on top.

The seabed dig surfaced a stronger result than expected: the exact
2D partition closed form is right there in the same file as A2's
factor theorem.\footnote{In \texttt{QuarticSextic.lean}, named
\texttt{partitionFunction\_quarticSextic\_eq}, sibling to
\texttt{partitionFunction\_quarticSextic\_factor} and T1's moment
factor theorem.} So Z2 is literally Z1 at $(j, k) = (0, 0)$ with the
seabed-provided exact identity replacing Z1's
separability-composition step. Even cleaner.

\section{The thing formalised}

Three theorems in \texttt{Laplace.TwoD}:

\begin{enumerate}[leftmargin=1.5em]
\item \textbf{Exact reformulation.} For $t > 0$,
  $Z_{2D}(t) = C''\cdot t^{-1/4 - 1/6}$ with
  $C'' = \tfrac{1}{6} \cdot 24^{1/4} \cdot 720^{1/6}
  \cdot \Gamma(1/4)\cdot \Gamma(1/6)$.
\item \textbf{Rescaled \texttt{Tendsto}.} As $t \to \infty$,
  $t^{1/4 + 1/6}\cdot Z_{2D}(t) \to C''$.
\item \textbf{\texttt{IsEquivalent}.} At $\mathrm{atTop}$,
  $Z_{2D}(t) \;\sim\; C''\cdot t^{-1/4 - 1/6}$.
\end{enumerate}

Named constant: \texttt{quarticSexticPartitionConst}. File:
sibling to V2's and Z1's asymptotic files, in the same
\texttt{Laplace/TwoD/} folder.\footnote{Specifically
\texttt{QuarticSexticPartitionAsymptotic.lean}, alongside
\texttt{QuarticSexticMomentAsymptotic.lean} (V2) and
\texttt{QuarticSexticNumeratorAsymptotic.lean} (Z1).}

\section{Proof strategy}

For the exact reformulation: one \texttt{rw} of the seabed's
existing 2D-partition closed form (giving
$(1/2)(24/t)^{1/4}\Gamma(1/4) \cdot (1/3)(720/t)^{1/6}\Gamma(1/6)$),
then the same \texttt{Real.rpow} algebra V2 and Z1 used:
\texttt{div\_rpow} splits $(24/t)^{1/4}$, \texttt{rpow\_add}
combines $t^{-1/4}\cdot t^{-1/6} = t^{-1/4 - 1/6}$,
\texttt{rpow\_neg} bridges the negative exponents, and \texttt{ring}
closes. The two asymptotic forms follow by the same
eventual-equality idiom.

\section{Roadblocks and resolutions}

None. Z2 is the simplest of today's three mirror-shaped tides: a
single seabed identity rewrite plus the standard \texttt{rpow}
collapse. No \texttt{field\_simp} pitfall this time (lesson from V2
carried over). \texttt{lake build} green on first compile.

\section{What was Mathlib, what was new}

\paragraph{From Mathlib.} The same set of API V2/Z1 used: the
\texttt{Real.rpow} algebra suite (\texttt{div\_rpow},
\texttt{rpow\_add}, \texttt{rpow\_neg}), the
\texttt{Tendsto}-to-constant idiom, and \texttt{IsEquivalent}'s
\texttt{refl}+\texttt{congr\_left} pair.

\paragraph{From the seabed.} One lemma: the existing exact 2D
partition closed form.\footnote{\texttt{partitionFunction\_quarticSextic\_eq}
in \texttt{Laplace/TwoD/QuarticSextic.lean}.} Nothing else.

\paragraph{New here.} Three theorems plus the named constant
\texttt{quarticSexticPartitionConst}. The architectural value is in
completing the V2/Z1/Z2 triple --- moment / numerator / partition
all now have parallel asymptotic packagings, in three sibling files.

\section{Lessons}

\paragraph{Three is the right cadence for mirror-shaped tides.} V2
established the packaging shape against a freshly-deliberated GPT
consult; Z1 mirrored it with one structural change (the seabed dig
that found the existing factor lemma); Z2 mirrored both at minimal
cost. The pattern is reproducible: once a packaging shape has been
deliberated for one object, applying it to closely related objects
is cheap and skip-GPT-safe. Worth recording as a tide-template
recipe.

\paragraph{The 2D quartic-sextic story has a natural ``triple''
shape.} Moment, numerator, partition are the three objects to which
the asymptotic-packaging recipe applies; the partition is the
``unit'' object ($j = k = 0$ in the numerator family), so Z2
specialises Z1. Future programmes on other separable 2D potentials
should be able to ship the same three-theorem-times-three structure
with minimal additional deliberation.

\section{Follow-ups}

\begin{itemize}[leftmargin=1.5em]
\item \emph{Generic-$k$ 1D and 2D analogues.} The whole moment /
  numerator / partition triple should extend to pure $x^{2k}/(2k)!$
  potentials in 1D (using N1's J-function machinery) and to the
  separable 2D combinations of any pair of such 1D potentials. Out
  of scope for a single Tide; this would be a programme that the
  V2/Z1/Z2 triple has effectively prototyped.
\item \emph{Promote the three named constants to a single shared
  definition.} \texttt{quarticSexticMomentConst},
  \texttt{quarticSexticNumeratorConst}, and
  \texttt{quarticSexticPartitionConst} share structure; a single
  parametrised definition might consolidate them. Currently each is
  used once; defer until a fourth consumer appears.
\item \emph{\texttt{\textbackslash leanref}-tag the primer / grammar
  paper} with the full V2/Z1/Z2 triple. Pairs with the existing
  paper-tagging backlog.
\end{itemize}

\section*{Acknowledgements}

No GPT consult was used for Z2 (third mirror-shaped tide of the
day). The deliberation work inherited from V2's consult, which set
the three-theorem packaging shape and the \texttt{tendsto\_congr'}
idiom, and applied unchanged here. No numerical check was needed
beyond inheriting Z1's $(j, k) = (0, 0)$ row, which verified the
closed form to relative error ${\sim}10^{-16}$.

\end{document}
